\documentclass[11pt]{article}
\usepackage[margin=1in]{geometry}
\usepackage{amsmath,amssymb,amsthm}
\newtheorem{theorem}{Theorem}
\newtheorem{lemma}{Lemma}
\newcommand{\Fp}{\mathbb F_p}
\title{Uniform atoms for nearconstant inputs in the configuration model}
\author{Research working note; proof supplied, priority unverified}
\date{November 15, 2026}
\begin{document}
\maketitle

Fix $d\ge3$ and a finite field $\Fp$ with $p$ elements. A uniform pairing
of $dN$ labelled stubs, $d$ at each of $N$ vertices, defines the regular
configuration-model adjacency matrix $C$. Loops contribute two on the
diagonal, so $C\mathbf1=d\mathbf1$. Assume $dN$ is even. In particular,
odd $N$ is permitted when $d$ is even. No assertion about shifted
invertibility follows merely from unshifted invertibility; the atom
estimate here is uniform over every target vector.

\begin{theorem}\label{thm:nearconstant}
There exists $\eta=\eta(d,p)>0$ such that
\[
 \sum_{\substack{q\in\Fp^N\ \mathrm{nonconstant}\\
 \min_{a\in\Fp}|\{v:q_v\ne a\}|\le\eta N}}
 \sup_{r\in\Fp^N}\Pr(Cq=r)
 =O_{d,p}(N^{1-d/2}).
\]
Thus the sum tends to zero along every admissible parity subsequence.
The same upper bound holds when $C$ is conditioned on simplicity, at
the cost of the usual fixed-$d$ bounded conditioning factor.
\end{theorem}

The proof uses only direct stub exposure. It neither assumes a random
target nor imports a mixing theorem for odd vertex count.

\begin{lemma}[An atom bound with arbitrary output]\label{lem:atom}
Put $K_d=2(2d+1)$. If $z\in\Fp^N$ has exactly $k$ nonzero coordinates,
$1\le k\le N/4$, then, provided $K_dk/N\le1$,
\[
 \sup_r\Pr(Cz=r)
 \le2^{dk+1}\left(\frac{K_dk}{N}\right)^{dk/2}.
\]
\end{lemma}
\begin{proof}
Fix a target $r$, and let $S=\operatorname{supp}z$, so $|S|=k$.
Let $T=\{v\notin S:r_v\ne0\}$. If $|T|>dk$, the probability is zero:
every such vertex must have at least one stub paired to a stub of $S$.
Otherwise $|T|\le dk$.

Expose partners of stubs of $S$, always taking the first still-unpaired
stub of $S$ in a fixed order. The conditional partner is uniform over
all remaining stubs other than the current one. An exposure is called
constrained if its partner vertex belongs to $S\cup T$ or is an outside
vertex previously hit in this exposure process. All other exposures
are fresh. There are at most $dk$ exposure steps. At every step the
number of eligible partner stubs at constrained vertices is at most
\[
 d\bigl(|S|+|T|+dk\bigr)\le d(2d+1)k.
\]
The total number of available partner stubs is at least
$dN-2dk+1\ge dN/2$, since $k\le N/4$. Therefore, whatever the past
history, the conditional probability that a next step exists and is
constrained is at most $K_dk/N$.

On the event $Cz=r$, any vertex outside $S\cup T$ hit by a stub from
$S$ must be hit at least twice: exactly one hit would contribute one
nonzero coordinate of $z$ and could not sum to its required output zero.
This remains valid with repeated graph edges and in every characteristic.
Let $I$ count exposed pairs with both stubs in $S$, $J$ exposed edges
from $S$ to $T$, $F$ fresh exposures outside $S\cup T$, and $R$ repeat
exposures outside $S\cup T$. Then
\[
 dk=2I+J+F+R,\qquad F\le R.
\]
The number of constrained steps is $I+J+R$, and hence
\[
 2(I+J+R)-dk=J+R-F\ge0.
\]
There are at least $dk/2$ constrained steps.

Encode the realized process by its finite binary word of constrained
and fresh steps. Its length is at most $dk$, so fewer than
$2^{dk+1}$ words are possible. For any specified word with $c$
constrained steps, its probability is at most $(K_dk/N)^c$, by the
conditional bound above; specifying the other steps and the stopping
condition can only decrease that probability. On $Cz=r$, $c\ge dk/2$.
Summing over those words proves the claim.
\end{proof}

\begin{proof}[Proof of Theorem~\ref{thm:nearconstant}]
If $q=a\mathbf1+z$ with support size $k$, then
$Cq=da\mathbf1+Cz$. Thus subtracting a constant does not change the
largest atom. There are at most $p\binom Nk(p-1)^k$ such inputs for
each $k$; overcounting is harmless. Put
\[
 \beta=d/2-1>0,\qquad
 A_{d,p}=2^d e(p-1)K_d^{d/2}.
\]
Choose $\eta>0$ satisfying
\[
 \eta\le\min\left\{\frac14,\frac1{K_d},
            (2A_{d,p}e^\beta)^{-1/\beta}\right\}.
\]
By Lemma~\ref{lem:atom} and $\binom Nk\le(eN/k)^k$, the desired sum
is at most
\[
 2p\sum_{k=1}^{\lfloor\eta N\rfloor}
 \left[A_{d,p}(k/N)^\beta\right]^k.
\]
Writing $t_k=[A_{d,p}(k/N)^\beta]^k$, for $k+1\le\eta N$ we have
\[
 \frac{t_{k+1}}{t_k}
 =A_{d,p}\left(\frac{k+1}{N}\right)^\beta
        (1+1/k)^{\beta k}
 \le A_{d,p}\eta^\beta e^\beta\le\frac12.
\]
Consequently the sum is at most $4p A_{d,p}N^{-\beta}$, proving the
bound. For fixed $d$, the configuration model has simplicity probability
bounded below away from zero; dividing every atom bound by that
probability gives the conditional version.
\end{proof}

\paragraph{Sharp order in the unconditioned even-degree model.}
If $d$ is even, fix a vertex $v$ and nonzero $b\in\Fp$. The event that
all $d$ stubs of $v$ pair among themselves has probability
\[
 \frac{(d-1)!!}{(dN-1)(dN-3)\cdots(dN-d+1)}
 =\Theta_d(N^{-d/2}).
\]
On that event, $C(be_v)=db\,e_v$. Summing these lower bounds over
the $N(p-1)$ distinct inputs $be_v$ proves a lower bound of order
$N^{1-d/2}$ for the nearconstant atom sum, once $\eta N\ge1$.
This order-sharpness example uses loops and does not survive
conditioning on simplicity.

\paragraph{Scope for RSBM.}
The estimate proves the nearconstant part of a possible finite-field
mixing theorem for the configuration model, uniformly over target
vectors. It does not supply the remaining balanced-input local central
limit estimate or the entropy bound for intermediate input types.
Those require separate arguments before concluding odd-size planted
eigenvalue simplicity.
\end{document}
